\documentclass[a4paper,11pt]{article}
\usepackage[margin=2.2cm]{geometry}
\usepackage{amsmath,amssymb}
\usepackage{booktabs}
\usepackage{array}
\usepackage{longtable}
\usepackage{tabularx}
\usepackage{makecell}
\usepackage{rotating}
\usepackage{hyperref}
\usepackage{cleveref}

\newtheorem{definition}{Definition}
\DeclareMathOperator{\polylog}{polylog}

\begin{document}

\title{Computational Complexity Comparison of FL Baselines}
\author{}
\date{}
\maketitle

\section{Notation}
\label{sec:notation}

\begin{table}[h]
\centering
\caption{Complexity-analysis notation}
\label{tab:comp-notation}
\renewcommand{\arraystretch}{1.15}
\small
\begin{tabular}{c>{\raggedright\arraybackslash}p{\dimexpr\linewidth-3.5cm}}
\toprule
\textbf{Symbol} & \textbf{Meaning} \\
\midrule
$d$ & Total number of model parameters \\
$N$ & Total number of clients \\
$K$ & Number of clients selected per round ($K\le N$) \\
$T$ & Total number of global rounds \\
$E$ & Number of local training epochs per round \\
$|\mathcal{D}_i|$ & Size of client $i$'s local dataset \\
$n$ & CKKS/RLWE polynomial modulus degree; slot count is $n/2$ \\
$J$ & Number of packs/ciphertexts for the full model, $J=\lceil 2d/n\rceil$ for merged CKKS packing \\
$k$ & Number of consensus packs selected by voting ($k\le J$) \\
$s$ & Per-client sparsity ratio in FedPHE, $0<s\le 1$ \\
$R$ & Prefetch window in ECPFL (rounds amortized into one decryption) \\
$L$ & Number of RNS moduli (levels) in multi-key CKKS ciphertexts \\
$\mathcal{O}_{\text{ntt}}$ & Cost of one NTT/INTT, $\mathcal{O}_{\text{ntt}}=\mathcal{O}(n\log n)$ \\
$\mathcal{O}_{\text{he-add}}$ & Cost of one ciphertext addition, $\mathcal{O}_{\text{he-add}}=\mathcal{O}(n)$ \\
$\mathcal{O}_{\text{he-smult}}$ & Cost of one ciphertext scalar multiplication, $\mathcal{O}_{\text{he-smult}}=\mathcal{O}(n)$ \\
\bottomrule
\end{tabular}
\end{table}

All complexities are stated \textbf{per global round} unless otherwise noted. The local-training cost $\mathcal{C}_{\text{train}}=O(E\,|\mathcal{D}_i|\,d)$ is identical for every scheme and is therefore omitted from the dominant cryptographic comparisons, but it is included for completeness in \cref{sec:per-scheme}.

\section{Complexity Summary}
\label{sec:summary}

\begin{sidewaystable}[h]
\centering
\caption{Per-round computational complexity comparison}
\label{tab:complexity-summary}
\renewcommand{\arraystretch}{1.25}
\small
\begin{tabularx}{\textheight}{|l|X|X|X|}
\hline
\textbf{Scheme} & \textbf{Client (per selected client)} & \textbf{Server} & \textbf{KGC / Authority} \\
\hline
FedAvg+CKKS & $\mathcal{O}(d\log n)$ enc/dec & $\mathcal{O}(Kd)$ HE ops $+\mathcal{O}(d\log n)$ dec & One-time only \\
\hline
MASER & $\mathcal{O}(d\log n + d_w)$ enc/partial-dec/mask & $\mathcal{O}(K^2 d L)$ agg $+\mathcal{O}(KdL)$ merge & $\mathcal{O}(NdL)$ one-time pk aggregation \\
\hline
FedPHE & $\mathcal{O}(sd\log n)$ enc $+\mathcal{O}(d\log n)$ dec & $\mathcal{O}(Ksd\log n)$ HE mult/add & None in reference impl. \\
\hline
ESAFL & $\mathcal{O}(d\log n)$ enc $+\mathcal{O}(d\log n)$ dec & $\mathcal{O}(Kd)$ HE additions & One-time only \\
\hline
ECPFL (proposed) & 
\makecell[l]{$\mathcal{O}(d+NJ)$ vote+perturbation \\ $+\mathcal{O}(k n\log n)$ enc \\ amortized $\mathcal{O}(\frac{k n\log n}{R})$ prefetch dec}
&
\makecell[l]{$\mathcal{O}(KJ)$ vote aggregation \\ $+\mathcal{O}(J\log J)$ consensus \\ $+\mathcal{O}(Kk)$ HE aggregation \\ $+\mathcal{O}(k)$ prefetch delta}
&
\makecell[l]{$\mathcal{O}(NKJ)$ compensatory perturbation \\ $+\mathcal{O}((N-K)k n\log n)$ \\ compensatory ciphertexts}
\\
\hline
\end{tabularx}
\end{sidewaystable}

\noindent
The bottom-line comparison is:
\begin{itemize}
    \item \textbf{FedAvg+CKKS} and \textbf{ESAFL} have the lowest server aggregation cost, $\mathcal{O}(Kd)$, because ciphertexts do not grow during aggregation.
    \item \textbf{MASER} pays a multiplicative $K$ penalty on the server ($\mathcal{O}(K^2 dL)$) because its multi-key ciphertexts grow from 2 to $K+1$ components when homomorphically summed.
    \item \textbf{FedPHE} keeps server aggregation at $\mathcal{O}(Ksd\log n)$, but the $\log n$ factor appears because it uses CKKS ciphertext-plaintext multiplications for weighted aggregation.
    \item \textbf{ECPFL} adds voting/perturbation overhead $\mathcal{O}(NJ)$ per active client and $\mathcal{O}(NKJ)$ at the KGC, but reduces the encryption/upload volume to $k\le J$ packs via consensus sparsity and amortizes decryption over $R$ rounds via prefetching.
\end{itemize}

\section{Detailed Derivations}
\label{sec:per-scheme}

For each scheme we derive the complexity from its core design logic, not from simulation-code shortcuts.

\subsection{FedAvg+CKKS}
\label{ssec:fedavg-ckks}

\textbf{Design.} Standard FedAvg with a single CKKS key pair. All selected clients encrypt their local updates; the server aggregates ciphertexts homomorphically and decrypts the global model for evaluation.

Let $J=\lceil 2d/n\rceil$ be the number of ciphertexts needed for the full model under merged slot packing.

\paragraph{Client side.} Each selected client
\begin{enumerate}
    \item downloads $J$ ciphertexts and decrypts them: $J$ CKKS decryptions;
    \item trains locally and computes the model delta;
    \item encrypts the delta into $J$ ciphertexts: $J$ CKKS encryptions.
\end{enumerate}
Each CKKS encryption or decryption costs $\mathcal{O}_{\text{ntt}}=\mathcal{O}(n\log n)$. Therefore
\begin{equation}
\boxed{T_{\text{client}}^{\text{FedAvg+CKKS}} = 2J\cdot\mathcal{O}(n\log n) + \mathcal{O}(d) = \mathcal{O}(d\log n).}
\end{equation}

\paragraph{Server side.} The server receives $K\cdot J$ ciphertexts. For each of the $J$ packs it computes a weighted average over $K$ ciphertexts, which needs $K$ scalar multiplications and $K-1$ additions. Scalar mult and addition are both $\mathcal{O}(n)$. The server also decrypts $J$ ciphertexts for evaluation. Hence
\begin{equation}
\boxed{T_{\text{server}}^{\text{FedAvg+CKKS}} = J\cdot K\cdot\mathcal{O}(n) + J\cdot\mathcal{O}(n\log n) = \mathcal{O}(Kd) + \mathcal{O}(d\log n).}
\end{equation}

\paragraph{KGC.} None; the single secret key is generated once and shared with the server.

\subsection{MASER}
\label{ssec:maser}

\textbf{Design.} MASER uses multi-key CKKS (xMK-CKKS). Each client keeps its own secret key $sk_i$; a Key Manager aggregates the public-key components into a joint public key $pk$. Clients encrypt only parameters selected by a global majority-vote mask. The server sums ciphertexts and clients collectively perform partial decryption.

Let $m$ be the number of parameters selected by the global mask and $J=\lceil m/(n/2)\rceil$ the number of slices/ciphertexts. Let $d_w$ be the number of weight parameters used in mask voting.

\paragraph{Client side (per effective global round).} MASER uses three sub-rounds (mask, encrypt, decrypt) per global update. Each client performs:
\begin{enumerate}
    \item local training and magnitude-based local mask generation: $\mathcal{O}(d_w)$;
    \item encrypt $J$ slices under the aggregated public key: $J$ CKKS encryptions at $\mathcal{O}(n\log n + nL)$ each;
    \item partial-decrypt $J$ aggregated multi-key ciphertexts: $J$ partial decryptions at $\mathcal{O}(nL)$ each.
\end{enumerate}
Thus
\begin{equation}
\boxed{T_{\text{client}}^{\text{MASER}} = J\cdot\mathcal{O}(n\log n + nL) + \mathcal{O}(d_w) = \mathcal{O}(d\log n + d_w).}
\end{equation}

\paragraph{Server side (per effective global round).} The server
\begin{enumerate}
    \item aggregates $K$ local masks: $\mathcal{O}(K d_w)$;
    \item homomorphically sums $K$ ciphertexts for each of $J$ slices. Because xMK-CKKS ciphertexts grow from 2 components to $K+1$ components, the cost of summing $K$ ciphertexts is $\mathcal{O}(K^2 n L)$ per slice;
    \item merges $K$ partial-decryption shares and decodes each of $J$ slices: $\mathcal{O}(K n L + n\log n)$ per slice.
\end{enumerate}
Therefore
\begin{equation}
\boxed{T_{\text{server}}^{\text{MASER}} = J\cdot\mathcal{O}(K^2 nL) + J\cdot\mathcal{O}(KnL + n\log n) + \mathcal{O}(Kd_w) = \mathcal{O}(K^2 dL + KdL + d\log n + Kd_w).}
\end{equation}
The dominant term is $\mathcal{O}(K^2 dL)$.

\paragraph{KGC / Key Manager.} One-time aggregation of $N$ public keys: $\mathcal{O}(N n L)$.

\subsection{FedPHE}
\label{ssec:fedphe}

\textbf{Design.} FedPHE uses packed single-key CKKS (optionally Paillier/BFV). Each client sparsifies its update to an $s$ fraction of packs, encrypts the selected packs, and uploads a binary mask. The server performs ciphertext-plaintext multiplication for contribution-aware weighted aggregation. Decryption is done at the clients.

Let $J=\lceil 2d/n\rceil$ and let each client upload $sJ$ packs in expectation.

\paragraph{Client side.} Each selected client
\begin{enumerate}
    \item trains locally;
    \item scores and selects top-$sJ$ packs: $\mathcal{O}(d + J\log J)$;
    \item encrypts $sJ$ packs: $sJ\cdot\mathcal{O}(n\log n)$;
    \item downloads and decrypts the global model in $J$ packs: $J\cdot\mathcal{O}(n\log n)$;
    \item fills zero positions from local parameters: $\mathcal{O}(d)$.
\end{enumerate}
Hence
\begin{equation}
\boxed{T_{\text{client}}^{\text{FedPHE}} = sJ\cdot\mathcal{O}(n\log n) + J\cdot\mathcal{O}(n\log n) + \mathcal{O}(d) = \mathcal{O}(sd\log n + d\log n) = \mathcal{O}(d\log n).}
\end{equation}

\paragraph{Server side.} The server
\begin{enumerate}
    \item aggregates $K$ masks of length $J$: $\mathcal{O}(KJ)$;
    \item for each active pack, performs a CKKS ciphertext-plaintext multiplication and addition. There are $sJ$ active packs per client and $K$ clients, so $KsJ$ mult/add operations at $\mathcal{O}(n\log n)$ each.
\end{enumerate}
Thus
\begin{equation}
\boxed{T_{\text{server}}^{\text{FedPHE}} = KsJ\cdot\mathcal{O}(n\log n) + \mathcal{O}(KJ) = \mathcal{O}(Ksd\log n).}
\end{equation}

\paragraph{Optional client selection.} MinHash-based selection costs $\mathcal{O}(Ld)$ per client and $\mathcal{O}(N^2L)$ server clustering, but it is typically run every several rounds, not every round.

\paragraph{KGC.} None in the reference implementation; the secret key is shared with all clients.

\subsection{ESAFL}
\label{ssec:esafl}

\textbf{Design.} ESAFL uses a multi-client RLWE additively homomorphic scheme (ESHE). Each client has its own secret key $s_i$; the global decryption key is $S=\sum_i s_i$. A per-round public polynomial $a^{(t)}$ is generated locally by all parties from a common seed via PRPG. Ciphertexts are aggregated by simple addition, and decryption is one-step using $S$.

Let $J=\lceil 2d/n\rceil$.

\paragraph{Client side.} Each selected client
\begin{enumerate}
    \item trains locally;
    \item computes PRPG output $a^{(t)}$: $\mathcal{O}(n)$;
    \item encodes and encrypts the model update into $J$ ciphertexts: $J\cdot\mathcal{O}(n\log n)$;
    \item downloads and decrypts the aggregated ciphertext with $S$: $J\cdot\mathcal{O}(n\log n)$.
\end{enumerate}
Therefore
\begin{equation}
\boxed{T_{\text{client}}^{\text{ESAFL}} = 2J\cdot\mathcal{O}(n\log n) + \mathcal{O}(n) + \mathcal{O}(d) = \mathcal{O}(d\log n).}
\end{equation}

\paragraph{Server side.} The server receives $K\cdot J$ ciphertexts and homomorphically sums them. Each addition is $\mathcal{O}(n)$. Hence
\begin{equation}
\boxed{T_{\text{server}}^{\text{ESAFL}} = KJ\cdot\mathcal{O}(n) = \mathcal{O}(Kd).}
\end{equation}

\paragraph{KGC.} One-time setup: generate $N$ keys, compute $S=\sum_i s_i$, distribute seeds. No per-round KGC work.

\subsection{ECPFL (Proposed)}
\label{ssec:ecpfl}

\textbf{Design.} ECPFL builds on ESAFL and adds (i) pack-level voting with a global consensus set $\mathcal{K}^{(t)}$, (ii) pairwise perturbation for privacy-preserving vote aggregation, and (iii) ciphertext prefetching so inactive clients amortize decryption over $R$ rounds.

Let $k=|\mathcal{K}^{(t)}|$ be the number of consensus packs, with $k\le J$.

\subsubsection{Active client}
\label{sssec:ecpfl-active}

An active client performs:
\begin{enumerate}
    \item local training: $\mathcal{O}(E|\mathcal{D}_i|d)$;
    \item residual compensation $\tilde{\mathbf{g}}_i=\mathbf{g}_i+\mathbf{r}_i$: $\mathcal{O}(d)$;
    \item score all $J$ packs: $\mathcal{O}(d)$;
    \item select local Top-$k$ packs: $\mathcal{O}(J)$ with linear selection or $\mathcal{O}(J\log J)$ with sorting;
    \item build raw vote bitmap $\mathbf{v}_i\in\{0,1\}^J$: $\mathcal{O}(J)$;
    \item generate pairwise perturbation $\mathbf{z}_i^{(t)}$: $N-1$ PRG calls each of length $J$, costing $\mathcal{O}(NJ)$;
    \item obfuscate vote $\hat{\mathbf{v}}_i=[\mathbf{v}_i+\mathbf{z}_i]_M$: $\mathcal{O}(J)$;
    \item upload obfuscated vote: $\mathcal{O}(J)$ communication;
    \item receive consensus set $\mathcal{K}^{(t)}$: $\mathcal{O}(J)$ or $\mathcal{O}(k)$ communication;
    \item generate PRPG $a^{(t)}$: $\mathcal{O}(n)$;
    \item encrypt $k$ consensus packs: $k\cdot\mathcal{O}(n\log n)$;
    \item update residual $\mathbf{r}_i^{(t)}$: $\mathcal{O}(d)$.
\end{enumerate}
Adding the non-training terms gives
\begin{equation}
\boxed{T_{\text{client, active}}^{\text{ECPFL}} = \mathcal{O}(d+NJ) + k\cdot\mathcal{O}(n\log n) + \mathcal{O}(E|\mathcal{D}_i|d).}
\end{equation}

\subsubsection{Inactive / prefetch client}
\label{sssec:ecpfl-inactive}

An inactive client participating in prefetching:
\begin{enumerate}
    \item receives and stores/merges prefetch deltas for $R$ rounds: $\mathcal{O}(Rk)$ ciphertext references and $\mathcal{O}(Rk)$ homomorphic additions if accumulation is implemented;
    \item when selected, performs one CKKS decryption of the accumulated $R$-round composite ciphertext: $k\cdot\mathcal{O}(n\log n)$, amortized over $R$ rounds.
\end{enumerate}
Thus the amortized per-round prefetch cost is
\begin{equation}
\boxed{T_{\text{client, prefetch}}^{\text{ECPFL}} = \mathcal{O}(k) + \frac{k\cdot\mathcal{O}(n\log n)}{R} = \mathcal{O}\!\left(k + \frac{k n\log n}{R}\right).}
\end{equation}

\subsubsection{Server}
\label{sssec:ecpfl-server}

The server performs:
\begin{enumerate}
    \item client selection / scheduling (e.g., sticky or sketch-based): $\mathcal{O}(N\log N)$ or $\mathcal{O}(N)$, typically not on the critical training path;
    \item receive $K$ obfuscated votes of length $J$: $\mathcal{O}(KJ)$ communication;
    \item ESP-Consensus: sum obfuscated votes and add KGC compensatory perturbation, $\mathcal{O}(KJ)$ modular additions; then compute Top-$k$ consensus packs, $\mathcal{O}(J)$ or $\mathcal{O}(J\log J)$;
    \item broadcast consensus set: $\mathcal{O}(J)$ or $\mathcal{O}(k)$ communication;
    \item receive $K\cdot k$ ciphertexts;
    \item homomorphically aggregate per consensus pack over $K$ clients: $Kk$ additions/scalar mults at $\mathcal{O}(n)$ each;
    \item update encrypted global state: $\mathcal{O}(k)$;
    \item compute prefetch delta. In the ESAFL path the aggregated ciphertext is reused as the delta: $\mathcal{O}(k)$. In the standard path the server decrypts the global state, computes a plaintext difference, and re-encrypts: $\mathcal{O}(J n\log n)$;
    \item queue update and per-client fetch preparation: $\mathcal{O}(k)$ per client needing a merge.
\end{enumerate}
Ignoring the expensive standard-path fallback, the dominant server cost is
\begin{equation}
\boxed{T_{\text{server}}^{\text{ECPFL}} = \mathcal{O}(KJ) + \mathcal{O}(J\log J) + \mathcal{O}(Kk) + \mathcal{O}(k) = \mathcal{O}(KJ + J\log J + Kk).}
\end{equation}

\subsubsection{KGC}
\label{sssec:ecpfl-kgc}

The KGC performs one-time setup and per-round compensation:
\begin{enumerate}
    \item one-time: generate $N$ keys, $S=\sum_i s_i$, PRPG seed $B$, and $\binom{N}{2}$ pairwise seeds: $\mathcal{O}(N^2)$;
    \item per round: receive inactive set $\bar{\mathcal{S}}^{(t)}$;
    \item compute compensatory perturbation $\mathbf{z}_{\mathrm{comp}}^{(t)}=\sum_{k\in\bar{\mathcal{S}}^{(t)}}\mathbf{z}_k^{(t)}$: $\mathcal{O}((N-K)NJ)$, or equivalently $\mathcal{O}(NKJ)$ using the active-client optimization $\mathbf{z}_{\mathrm{comp}}^{(t)}=-\sum_{i\in\mathcal{S}^{(t)}}\mathbf{z}_i^{(t)}$;
    \item generate compensatory ciphertexts for each inactive client and each consensus pack: $(N-K)k$ CKKS encryptions at $\mathcal{O}(n\log n)$ each;
    \item prefetch key-alignment masks: $\mathcal{O}(n\log n)$ per client per round.
\end{enumerate}
Hence
\begin{equation}
\boxed{T_{\text{KGC}}^{\text{ECPFL}} = \mathcal{O}(NKJ) + (N-K)k\cdot\mathcal{O}(n\log n) + N\cdot\mathcal{O}(n\log n).}
\end{equation}

\section{Unified Comparison}
\label{sec:unified}

The summary in \cref{tab:complexity-summary} expresses each scheme in its own notation, which makes direct comparison difficult. \cref{tab:unified-comparison} therefore uses a single set of metrics---dominant client cost, dominant server cost, KGC online cost, upload/download ciphertext counts, rounds per global update, dropout tolerance, and suitable deployment regime---to expose the trade-offs more clearly.

\begin{sidewaystable}[h]
\centering
\caption{Unified comparison of FL schemes (one effective global model update)}
\label{tab:unified-comparison}
\renewcommand{\arraystretch}{1.2}
\small
\begin{tabularx}{\textheight}{|l|>{\raggedright\arraybackslash}X|>{\raggedright\arraybackslash}X|>{\raggedright\arraybackslash}l|}
\hline
\textbf{Scheme} & \textbf{Client dominant cost} & \textbf{Server dominant cost} & \textbf{KGC online cost} \\
\hline
FedAvg+CKKS & $\mathcal{O}(d\log n)$ enc/dec & $\mathcal{O}(Kd)$ HE ops $+ \mathcal{O}(d\log n)$ dec & None \\
\hline
MASER & $\mathcal{O}(d\log n)$ enc/partial-dec & $\mathcal{O}(K^2 dL)$ agg $+ \mathcal{O}(KdL)$ merge & One-time pk aggregation \\
\hline
FedPHE & $\mathcal{O}(d\log n)$ enc/dec & $\mathcal{O}(Ksd\log n)$ HE mult/add & None \\
\hline
ESAFL & $\mathcal{O}(d\log n)$ enc/dec & $\mathcal{O}(Kd)$ HE additions & None \\
\hline
ECPFL & $\mathcal{O}(d+NJ)$ vote/perturbation $+ \mathcal{O}(k n\log n)$ enc & $\mathcal{O}(KJ + Kk)$ vote + HE agg & $\mathcal{O}(NKJ) + \mathcal{O}((N-K)k n\log n)$ \\
\hline
\end{tabularx}
\vspace{2mm}
\begin{tabularx}{\linewidth}{|l|>{\raggedright\arraybackslash}X|>{\raggedright\arraybackslash}X|>{\raggedright\arraybackslash}X|>{\raggedright\arraybackslash}X|}
\hline
\textbf{Scheme} & \textbf{Upload / Download} & \textbf{Rounds per update} & \textbf{Dropout robust} & \textbf{Regime} \\
\hline
FedAvg+CKKS & $J$ / $J$ & 1 & Limited & Cross-silo \\
\hline
MASER & $J$ / $J$ & 3 & No & Cross-silo \\
\hline
FedPHE & $sJ$ / $J$ & 1 & Yes & Cross-device \\
\hline
ESAFL & $J$ / $J$ & 1 & No & Cross-silo \\
\hline
ECPFL & $k$ / $k$ amortized & 1 & Partial & Cross-silo \\
\hline
\end{tabularx}
\end{sidewaystable}

\subsection{What the unified view reveals}
\label{ssec:unified-view}

\begin{itemize}
    \item \textbf{Server aggregation.} ESAFL and FedAvg+CKKS have the cheapest server aggregation, $\mathcal{O}(Kd)$. ECPFL is similar when $k\ll J$, while FedPHE pays an extra $\log n$ factor because it uses CKKS multiplication for weighting. MASER is the most expensive server-side scheme due to growing multi-key ciphertexts ($\mathcal{O}(K^2 dL)$).
    \item \textbf{Client upload volume.} ECPFL and FedPHE reduce upload relative to full-model schemes: ECPFL uploads $k$ consensus packs and FedPHE uploads $sJ$ packs. For typical $k\approx sJ$, both achieve comparable upload savings, but ECPFL also coordinates clients to encrypt the same $k$ packs, which simplifies server aggregation.
    \item \textbf{Client-side extra work.} ECPFL adds $\mathcal{O}(d+NJ)$ non-HE perturbation work per active client. This term dominates the client cost when $NJ \gg d\log n$, i.e., when the number of clients $N$ is large relative to the model dimension $d$.
    \item \textbf{KGC online work.} Only ECPFL requires a trusted third party to perform online work ($\mathcal{O}(NKJ)$ compensatory perturbation and compensatory ciphertext generation). MASER's KGC is one-time; FedAvg+CKKS, FedPHE, and ESAFL have no online KGC.
\end{itemize}

\subsection{Implications for cross-device FL}
\label{ssec:cross-device}

Cross-device FL is characterized by very large $N$ (thousands to millions), small selected fractions $K/N$, limited client compute/battery, intermittent connectivity, and cheap cloud servers but no per-round trusted third party. Viewed through this lens, ECPFL has both advantages and limitations:

\begin{itemize}
    \item \textbf{Advantage: reduced upload/download volume.} Consensus sparsity reduces the per-client upload from $J$ packs to $k$ packs, and prefetching amortizes the client decryption over $R$ rounds. Both are attractive for bandwidth- and battery-constrained devices.
    \item \textbf{Advantage: partial dropout tolerance.} Compensatory ciphertexts allow the server to aggregate over any subset of active clients without requiring all clients to be online, unlike MASER or ESAFL.
    \item \textbf{Limitation: pairwise perturbation does not scale to large $N$.} The $\mathcal{O}(NJ)$ client-side perturbation generation and the $\mathcal{O}(NKJ)$ KGC compensatory computation grow with $N$. For $N$ in the millions, this is prohibitive. The $\mathcal{O}(N^2)$ pairwise seed setup is also infeasible at cross-device scale.
    \item \textbf{Limitation: online KGC.} Cross-device FL typically avoids online trusted authorities. ECPFL's reliance on the KGC to compute and transmit $\mathbf{z}_{\mathrm{comp}}^{(t)}$ and compensatory ciphertexts every round is a deployment obstacle.
    \item \textbf{Conclusion.} In its current form ECPFL is better suited to the \textbf{cross-silo} regime (tens to hundreds of stable clients, moderate $N$, trusted KGC) than to pure \textbf{cross-device} FL. To make it cross-device friendly, the pairwise vote-perturbation mechanism would need to be replaced by a more scalable approach, such as small-group pairwise masking, a lightweight secure aggregator (SecAgg), or accepting that the server learns aggregate vote tallies without cryptographic masking.
\end{itemize}

\section{Discussion}
\label{sec:discussion}

\begin{itemize}
    \item \textbf{Server aggregation cost.} FedAvg+CKKS, ESAFL, and ECPFL enjoy ciphertext-size-preserving aggregation, so server HE work is linear in $K$ (or $Kk/J$ for ECPFL). MASER's multi-key ciphertexts grow with the number of contributors, producing a $\mathcal{O}(K^2)$ factor.
    \item \textbf{Client encryption cost.} ECPFL encrypts only $k$ consensus packs per client instead of $J$, giving $k/J$ savings on the dominant HE encryption term relative to ESAFL/FedAvg+CKKS. FedPHE achieves similar savings with per-client TopK sparsity but pays $\log n$ per aggregation operation because it uses CKKS multiplication for weighting.
    \item \textbf{Prefetching trade-off.} Prefetching adds $\mathcal{O}(k)$ server-side HE work per round. For clients, the decryption cost is $\mathcal{O}(k n\log n)$ but is paid only once every $R$ rounds, giving an amortized $\mathcal{O}(k n\log n/R)$ cost and avoiding a full per-round decryption.
    \item \textbf{Vote-privacy overhead.} The pairwise perturbation mechanism adds $\mathcal{O}(NJ)$ client computation and $\mathcal{O}(NKJ)$ KGC computation per round. This is the main asymptotic novelty of ECPFL and is acceptable when $N$ is moderate (cross-silo regime).
    \item \textbf{MASER vs. ECPFL.} Both use multi-key techniques, but MASER relies on partial decryption and growing ciphertexts, while ECPFL/ESAFL use a summed global key $S$ for one-step decryption. ECPFL further reduces per-client upload via consensus voting and amortizes download via prefetching.
\end{itemize}

\end{document}
